\begin{table}
\caption{Tabela 2. Correlação de Spearman, bruta e parcial (controles: latitude, água e altitude), entre a perturbação e a temperatura da estação seca pelo MODIS (superfície, diurna) e pelo ERA5 (máxima do ar a 2 m), entre células (perda antecedente) e ao longo do tempo (perda 2001–2012 contra a mudança entre as décadas), com IC 95% por blocos de 1°.}
\label{tab:08_gradiente::modis}
\begin{tabular}{lllrrrl}
\toprule
Bacia & Comparação & Fonte & Células & ρ bruto & ρ parcial & IC 95\% \\
\midrule
Xingu & entre células & MODIS & 668 & 0,86 & 0,81 & [+0,73; +0,87] \\
Xingu & entre células & ERA5 & 668 & 0 & -0,14 & [-0,28; +0,01] \\
Xingu & ao longo do tempo & MODIS & 668 & 0,72 & 0,72 & [+0,63; +0,79] \\
Xingu & ao longo do tempo & ERA5 & 668 & -0,15 & -0,18 & [-0,35; +0,01] \\
Tapajós & entre células & MODIS & 652 & 0,89 & 0,74 & [+0,66; +0,80] \\
Tapajós & entre células & ERA5 & 652 & -0,38 & -0,06 & [-0,25; +0,13] \\
Tapajós & ao longo do tempo & MODIS & 652 & 0,73 & 0,77 & [+0,68; +0,82] \\
Tapajós & ao longo do tempo & ERA5 & 652 & -0,07 & 0,04 & [-0,21; +0,25] \\
\bottomrule
\end{tabular}
\end{table}
